\documentclass[11pt,a4paper]{article}
\usepackage[T1]{fontenc}
\usepackage[utf8]{inputenc}
\PassOptionsToPackage{dvipsnames}{xcolor}
\usepackage{amsmath,amsthm,mathtools}
\usepackage{newpxtext,newpxmath}
\usepackage[margin=27mm,headheight=15pt]{geometry}
\usepackage{microtype}
\usepackage{booktabs,array,longtable}
\usepackage[dvipsnames]{xcolor}
\usepackage{enumitem}
\usepackage{fancyhdr}
\usepackage[most]{tcolorbox}
\usepackage{listings}
\usepackage{hyperref}
\usepackage[nameinlink,noabbrev]{cleveref}
\definecolor{Ink}{HTML}{15354B}
\definecolor{Accent}{HTML}{246878}
\definecolor{Pale}{HTML}{F1F6F8}
\hypersetup{colorlinks=true,linkcolor=Ink,citecolor=Accent,urlcolor=Accent,
  pdftitle={An Explicit Formula for the Minimal e-Monotone Infinitary Ordinal Sum},
  pdfauthor={Research note prepared with ChatGPT},
  pdfsubject={A proposed solution to Lipparini's Problem 6.2}}
\pagestyle{fancy}
\fancyhf{}
\fancyhead[L]{\small\textsc{The minimal $e$-monotone ordinal sum}}
\fancyhead[R]{\small Research note}
\fancyfoot[C]{\thepage}
\renewcommand{\headrulewidth}{0.35pt}
\setlength{\parskip}{3pt}
\setlength{\parindent}{1.3em}
\setcounter{tocdepth}{2}
\setlist{itemsep=2pt,topsep=4pt}
\newtheorem{theorem}{Theorem}[section]
\newtheorem{lemma}[theorem]{Lemma}
\newtheorem{proposition}[theorem]{Proposition}
\newtheorem{corollary}[theorem]{Corollary}
\theoremstyle{definition}
\newtheorem{definition}[theorem]{Definition}
\newtheorem{example}[theorem]{Example}
\newtheorem{remark}[theorem]{Remark}
\newcommand{\On}{\mathrm{On}}
\newcommand{\cf}{\operatorname{cf}}
\newcommand{\rk}{\operatorname{rk}}
\newcommand{\ns}{\mathbin{\#}}
\newcommand{\nprod}{\mathbin{\otimes}}
\newcommand{\bigns}{\mathop{\#}\displaylimits}
\newcommand{\OpN}{\mathcal N}
\newcommand{\OpS}{\mathcal S}
\newcommand{\FunA}{\mathsf A}
\newcommand{\FunB}{\mathsf B}
\newcommand{\FunC}{\mathsf C}
\newcommand{\FunU}{\mathsf U}
\newcommand{\dd}[2]{d(#1,#2)}
\newcommand{\tail}{\varepsilon}
\newcommand{\prof}{\operatorname{prof}}
\newcommand{\cseq}[1]{\mathbf c_{#1}}
\newcommand{\restr}{\mathbin{\upharpoonright}}
\newcommand{\emptyheads}{\varnothing}
\newcommand{\eord}{\mathrel{\prec_e}}
\newtcolorbox{resultbox}[1][]{colback=Pale,colframe=Accent,boxrule=0.5pt,
  arc=1pt,left=9pt,right=9pt,top=7pt,bottom=7pt,#1}
\lstset{basicstyle=\small\ttfamily,columns=fullflexible,keepspaces=true,
  frame=single,rulecolor=\color{Accent},breaklines=true,showstringspaces=false}

\begin{document}
\begin{titlepage}
\vspace*{10mm}
{\color{Accent}\large\textsc{Ordinal arithmetic \& order-theoretic ranks}\par}
\vspace{9mm}
{\color{Ink}\LARGE\bfseries An Explicit Formula for the Minimal\\[3pt]
\boldmath$e$-Monotone Infinitary Ordinal Sum\par}
\vspace{6mm}
{\large A proposed solution to Lipparini's Problem 6.2\par}
\vspace{9mm}
{Research note prepared with ChatGPT\par}
{19 September 2026\par}
\vspace{11mm}
\begin{abstract}
For a countable sequence of ordinals $x=(\alpha_i)_{i<\omega}$, let
$\tail(x)$ be the least threshold at or above which only finitely many entries
occur. Consider the least ordinal-valued operation that is weakly monotone
coordinatewise and strictly increases whenever an increased coordinate reaches
$\tail(x)$. We give a three-case, finite Cantor-normal-form formula for this
operation. It depends only on the threshold and the finite multiset of
exceptional entries at or above it.

The proof separates an arithmetic upper bound from a rank-theoretic lower
bound. The key upper-bound lemma controls changes of threshold; the lower
bound uses a published minimality theorem for Lipparini's weaker operation,
together with finite-product ranks. We also identify the operation as the
rank of countable ordinal multisets under injective domination. Its difference
from the weaker operation is classified exactly: a finite increment, a natural
summand $\omega$, or zero. Consequences include the height of the domination
poset on countable multisets with bounded entries and an exact finite-cap
analogue. A dependency-free Python implementation and 734,869 counted symbolic
and finite-poset checks accompany the article.
\end{abstract}
\vfill
\begin{resultbox}
\textbf{Research status.} This is an unrefereed proposed solution, with a complete
conventional argument relative to one explicitly cited published theorem.
It has not been independently refereed or formalized in a proof assistant.
The computations test the formula and its lemmas; they do not establish the
transfinite theorem or priority.
\end{resultbox}
\end{titlepage}

{\small\setlength{\parskip}{0pt}
\tableofcontents}
\newpage

\section{The problem and the principal formula}\label{sec:main}

\subsection{A specific open question}
Lipparini's 2026 article introduces a least infinitary operation with
strengthened monotonicity, denoted there by $\#^N$, but leaves its explicit
construction open. Problem 6.2 in the extended version dated 30 April 2026
asks for an explicit definition and a study of that operation
\cite[Problem 6.2; Theorem 2.6.a(A)]{Lipparini2026extended}.
We write it as $\OpN$. The journal counterpart is \cite{Lipparini2026published}.

The present note proposes such a definition. The task is not to manufacture
one monotone upper bound: many upper bounds are possible. The essential point
is to prove \emph{pointwise minimality}. Sections~\ref{sec:upper}
and~\ref{sec:lower} prove the two inequalities separately. The source input for
the lower inequality is isolated in \cref{prop:importedS}; the remainder of the
argument is supplied here.

The April 2026 version is the latest version located in the searches made for
this note. This establishes the provenance of the question, not the absence
of every possible later or unpublished solution. The present formula's
priority has not been independently certified.

\subsection{The threshold and the monotonicity requirement}
All ordinals are allowed; the sequences, but not necessarily their entries,
are countable. We work in ZFC. For $x=(\alpha_i)_{i<\omega}$ define
\begin{equation}\label{eq:cut}
 \tail(x)=\min\{e\in\On: \{i<\omega:\alpha_i\ge e\}\text{ is finite}\}.
\end{equation}
This is always a positive ordinal. Put
\begin{equation}\label{eq:exceptions}
 E(x)=\{i<\omega:\alpha_i\ge\tail(x)\}.
\end{equation}
The set of indices $E(x)$ is finite. We retain the \emph{multiplicities} of its
values: three exceptional entries equal to $a$ contribute three summands,
not one.

Write $x\le y$ for coordinatewise comparison. An operation $F$ is
\emph{admissible} when, for all $x=(\alpha_i)$ and $y=(\beta_i)$,
\begin{align}
 x\le y&\ \Longrightarrow\ F(x)\le F(y), \label{eq:weak}\\
 x\le y,\quad \exists j\bigl(\alpha_j<\beta_j\text{ and }
           \tail(x)\le\beta_j\bigr)
 &\ \Longrightarrow\ F(x)<F(y).\label{eq:strict}
\end{align}
The threshold in \eqref{eq:strict} belongs to the \emph{lower} sequence.
We use that formulation throughout the rank proof. The operation $\OpN$ is
the pointwise least admissible operation. These are Lipparini's conditions
(2.1) and (2.2e) \cite[Definition 2.4]{Lipparini2026extended}.

If $\tail(x)=z+1$, infinitely many entries of $x$ equal $z$. If
$\tail(x)=e$ is a limit, its nonexceptional entries are cofinal in $e$, and
$\cf(e)=\omega$. In both cases, finite changes to the sequence preserve the
threshold. These facts will also follow explicitly in \cref{sec:profiles}.

\subsection{Arithmetic used in the formula}\label{subsec:notation}
The symbol $+$ always means ordinary ordinal addition, $\ns$ means the
Hessenberg natural sum, and $\nprod$ means the natural product. Empty finite
natural sums are zero. For $a\le b$, define
\begin{equation}\label{eq:remainder}
 \dd{a}{b}=\text{the unique ordinal }d\text{ such that }a+d=b.
\end{equation}
This is an ordinary-addition remainder, \emph{not} natural subtraction.
For example, $\dd{\omega}{\omega^2}=\omega^2$.

Set
\begin{equation}\label{eq:A}
 \FunA(t)=t\nprod\omega.
\end{equation}
If $t=\omega^{\eta_1}c_1+\cdots+\omega^{\eta_r}c_r$ is in Cantor normal form,
then
\begin{equation}\label{eq:Acnf}
 \FunA(t)=\omega^{\eta_1+1}c_1+\cdots+\omega^{\eta_r+1}c_r.
\end{equation}
The successors $\eta_j+1$ in this expression are ordinary ordinal successors.
In particular, $\FunA(\omega+1)=\omega^2+\omega$, not $\omega^2$.

Write $t=\lambda+n$ uniquely, where $n<\omega$ and $\lambda$ is zero or a
nonzero limit. Define
\begin{equation}\label{eq:chi}
 \chi(t)=
 \begin{cases}
  1,&\lambda>0\text{ and the least CNF exponent of }\lambda
       \text{ is a successor},\\
  0,&\text{otherwise},
 \end{cases}
 \qquad
 \FunC(t)=\FunA(t)\ns\bigl(\omega\cdot\chi(t)\bigr).
\end{equation}
Thus $\chi(\omega+n)=\chi(\omega^2+n)=1$, whereas
$\chi(\omega^\omega+n)=0$ for finite $n$.

For a nonzero limit $e$, write
\begin{equation}\label{eq:lastmonomial}
 e=p+\omega^\xi,\qquad \xi>0,
\end{equation}
by deleting \emph{one copy} of the final monomial in its additive normal form.
For example, if $e=\omega^5+\omega^2\cdot3$, then
$p=\omega^5+\omega^2\cdot2$. We do not delete the whole final CNF term.

\subsection{The explicit result}
\begin{theorem}[Explicit formula]\label{thm:main}
Let $x=(\alpha_i)_{i<\omega}$, let $e=\tail(x)$, and put $E=E(x)$.
Then the least admissible operation is given by the following finite formulas.

If $e=z+1$ is a successor, then
\begin{equation}\label{eq:main-successor}
 \boxed{\displaystyle
 \OpN(x)=\FunC(z)\ns\bigns_{i\in E}\dd{z}{\alpha_i}.}
\end{equation}
If $e=p+\omega^\xi$ is a limit and $\xi$ is a successor, then
\begin{equation}\label{eq:main-limit-successor}
 \boxed{\displaystyle
 \OpN(x)=\FunA(e)\ns\bigns_{i\in E}
             \bigl(1+\dd{e}{\alpha_i}\bigr).}
\end{equation}
If $e=p+\omega^\xi$ is a limit and $\xi$ is a nonzero limit, then
\begin{equation}\label{eq:main-limit-limit}
 \boxed{\displaystyle
 \OpN(x)=\FunA(p)\ns\omega^\xi
                    \ns\bigns_{i\in E}\dd{p}{\alpha_i}.}
\end{equation}
The $1+d$ in \eqref{eq:main-limit-successor} is \emph{ordinary left addition}.
It is not $d+1$ or $1\ns d$.
\end{theorem}

This is a finite expression once the threshold and exceptional multiset have
been supplied. It contains no recurrence over predecessor sequences and no
minimization over competing operations. Computing the threshold of an
arbitrary program-defined infinite sequence is a different, generally
noncomputable task; see \cref{subsec:computability}.

\begin{example}[Two distinct kinds of limit behavior]\label[example]{ex:overview}
For a sequence cofinal below $\omega$ and having $m$ exceptional copies of
$\omega$, the value is $\omega^2+m$. The constant sequence with value $\omega$
has value $\omega^2+\omega$. In contrast, a sequence cofinal below
$\omega^\omega$ with $m$ exceptional copies of $\omega^\omega$ has value
$\omega^\omega\cdot(m+1)$, and the constant sequence whose entries are $\omega^\omega$ has
value $\omega^{\omega+1}$, without an extra $\omega$.
\end{example}

\section{Elementary ordinal arithmetic and finite-product ranks}\label{sec:arithmetic}

\subsection{Remainders, natural sums, and low-order terms}
We collect the arithmetic facts needed in the proof. Their explicit
formulation matters because ordinary and natural addition behave differently
at limits.

\begin{lemma}\label[lemma]{lem:arith}
For ordinals in the indicated domains:
\begin{enumerate}[label=(\roman*)]
\item If $a\le b$, the remainder $\dd{a}{b}$ exists uniquely, and the map
      $b\mapsto\dd{a}{b}$ is strictly increasing on $[a,\infty)$.
\item If $a\le b\le c$, then
      $\dd{a}{c}=\dd{a}{b}+\dd{b}{c}$.
\item Natural sum is associative, commutative, and strictly increasing in
      each argument. Also $a+b\le a\ns b$.
\item A finite natural sum of ordinals smaller than $\omega^\xi$ is again
      smaller than $\omega^\xi$, for every $\xi>0$.
\item If $\delta<\omega^\xi$, then
      $\FunA(\delta)<\omega^{\xi+1}$. If $\xi$ is a nonzero limit, the stronger
      inequality $\FunA(\delta)<\omega^\xi$ holds.
\item $\FunA(t+1)=\FunA(t)\ns\omega$. Every nonzero $\FunA(t)$, and every
      nonzero $\FunC(t)$, is a limit ordinal.
\end{enumerate}
\end{lemma}
\begin{proof}
The interval $[a,b)$ has a unique order type $d$, giving $a+d=b$;
strict increase of ordinary addition in its right argument gives uniqueness
and monotonicity. Splitting $[a,c)$ at $b$ proves (ii).

For (iii), align the finitely many exponents in the two Cantor normal forms
and add their coefficients. A strict comparison is preserved at the largest
exponent where the operands differ. Ordinary addition may erase terms of its
left argument, whereas natural addition retains them, proving the inequality.
For (iv), all exponents in the finitely many operands are below $\xi$, and
this remains so after collecting terms. Formula~\eqref{eq:Acnf} proves (v):
if $\eta<\xi$, then $\eta+1\le\xi$, and if $\xi$ is a limit, then
$\eta+1<\xi$. Finally, adding one changes just the finite coefficient of $t$;
after the shift of exponents this changes just the coefficient of $\omega$.
Every exponent in a nonzero $\FunA(t)$ or $\FunC(t)$ is positive.
\end{proof}

If $t$ is a nonzero limit, all exponents in $\FunA(t)$ are at least $2$.
Consequently
\begin{equation}\label{eq:Cfinite}
 \FunC(t+n)=\FunC(t)+\omega\cdot n\qquad(n<\omega).
\end{equation}
Also, if $L$ is a limit ordinal, $r<L$, and $n<\omega$, then $r+n<L$.
Indeed, a finite sequence of successors of an ordinal below a limit remains
below that limit.

\subsection{The recursion for natural sum}
\begin{lemma}\label[lemma]{lem:natural-recursion}
For any ordinals $a,b$,
\begin{equation}\label{eq:natural-recursion}
 a\ns b=\sup\left(
   \{(a'\ns b)+1:a'<a\}\ \cup\
   \{(a\ns b')+1:b'<b\}\right).
\end{equation}
An empty supremum is zero.
\end{lemma}
\begin{proof}
Strict monotonicity gives the upper bound. For the reverse inequality, take
$\gamma<a\ns b$ and align the CNFs of these ordinals. At the greatest
exponent $\xi$ where they differ, the coefficient in $a\ns b$ is larger.
At least one of $a,b$ has a positive coefficient there; suppose it is $a$.
Reduce that coefficient of $a$ by one, leave its higher terms unchanged,
and replace its lower terms by the part of $\gamma$ below $\omega^\xi$.
The resulting $a'$ is smaller than $a$. If the coefficient after reduction
still exceeds that of $\gamma$, then $a'\ns b>\gamma$; otherwise the inserted
lower terms ensure $a'\ns b\ge\gamma$. For $\xi=0$ the lower part is empty
and the same argument applies. Thus one of the terms on the right of
\eqref{eq:natural-recursion} exceeds $\gamma$.
\end{proof}

\begin{lemma}[Rank of a finite product]\label[lemma]{lem:product-rank}
Let $P_1,\ldots,P_h$ be well-founded posets, with ranks normalized by
$\rk(p)=\sup_{q<p}(\rk(q)+1)$. In the coordinatewise product, ordered strictly
when at least one coordinate is strictly smaller,
\begin{equation}\label{eq:product-rank}
 \rk(p_1,\ldots,p_h)=\rk(p_1)\ns\cdots\ns\rk(p_h).
\end{equation}
A map preserving strict order cannot decrease rank.
\end{lemma}
\begin{proof}
The last assertion follows immediately by induction on the source rank.
For two factors, induction and strict monotonicity of natural sum give the
upper bound in \eqref{eq:product-rank}. If $a'<\rk(p)$, some predecessor of
$p$ has rank at least $a'$: otherwise the defining supremum for $\rk(p)$
would be at most $a'$. Apply this observation in each coordinate and use
\cref{lem:natural-recursion} for the reverse inequality. Induction on the
number of factors proves the statement. Notice that this argument uses the
recursion for natural sum, not an incorrect assertion of separate continuity
of natural sum at every limit.
\end{proof}

\section{Finite profiles and the exact rank interpretation}\label{sec:profiles}

\subsection{Which tail thresholds can occur?}
\begin{lemma}\label[lemma]{lem:cutfacts}
For every countable ordinal sequence $x$:
\begin{enumerate}[label=(\roman*)]
\item $\tail(x)>0$, and finite modifications leave $\tail(x)$ unchanged.
\item If $\tail(x)=z+1$, infinitely many entries equal $z$.
\item If $\tail(x)=e$ is a limit, $\cf(e)=\omega$ and the nonexceptional
      entries are cofinal in $e$.
\item Conversely, every positive successor and every nonzero limit of
      cofinality $\omega$ occurs as a threshold, with any prescribed finite
      multiset of exceptional entries at or above it.
\end{enumerate}
\end{lemma}
\begin{proof}
There are infinitely many entries at or above zero. Whether the set of
entries at or above a fixed ordinal is finite is unaffected by finitely many
changes, proving (i). If $e=z+1$, infinitely many entries are at or above $z$
and only finitely many exceed it, proving (ii).

Suppose $e$ is a limit. Below each $\gamma<e$ there cannot be an eventual
bound smaller than $e$: specifically, there are infinitely many entries at
or above $\gamma$ and only finitely many at or above $e$. The entries below
$e$ are therefore cofinal. A countable cofinal subset of a nonzero limit
has cofinality $\omega$.

For (iv), use a constant $z$-tail for the cut $z+1$, and a countable cofinal
tail below $e$ for a limit cut. Append the prescribed exceptional entries.
Finite additions to a sequence do not change its threshold.
\end{proof}

Call the ordinals in (iv) \emph{realizable cuts}. The following finite datum
forgets the order of coordinates and all inessential details of the tail.

\begin{definition}\label[definition]{def:profile}
The profile of $x$ is
\[
 \prof(x)=(e;A),\qquad e=\tail(x),\quad
 A=\{\!\{\alpha_i:i\in E(x)\}\!\},
\]
where $A$ is a finite multiset. For profiles $(e;A)$ and $(f;B)$ set
\begin{equation}\label{eq:profileorder}
 (e;A)\preceq(f;B)
 \quad\Longleftrightarrow\quad
 e\le f\ \text{ and }\
 |\{a\in A:a\ge\gamma\}|\le|\{b\in B:b\ge\gamma\}|
 \quad\text{for every }\gamma\ge f.
\end{equation}
Multiset cardinalities include multiplicities. The strict part of this
partial order is denoted by $\prec$.
\end{definition}

For computation, take the entries of $A$ that are at least $f$ in decreasing
order, and compare them with the decreasing list $B$. There must be at least
as many entries in $B$, and each corresponding entry of $B$ must be at least
as large. This is equivalent to \eqref{eq:profileorder}.

\subsection{Countable multiset domination}
\begin{lemma}[Matching characterization]\label[lemma]{lem:matching}
There is an injection $h:\omega\to\omega$ such that
$\alpha_i\le\beta_{h(i)}$ for all $i$ if and only if
$\prof(x)\preceq\prof(y)$.
\end{lemma}
\begin{proof}
An injection implies the threshold comparison: if $\tail(x)>\tail(y)=f$,
infinitely many entries of $x$ are at least $f$, but only finitely many
entries of $y$ are. It also implies all the finite cardinality inequalities
in \eqref{eq:profileorder}.

Conversely, first match the finitely many entries of $x$ at or above $f$
into the exceptional multiset of $y$, using the decreasing-list criterion.
Every remaining entry $a$ of $x$ is below $f$. By the definition of $f$,
infinitely many entries of $y$ are at least $a$. Enumerate the remaining
indices of $x$ and match them greedily, avoiding the finitely many indices
already used at each step. There is always an available choice.
\end{proof}

\begin{lemma}\label[lemma]{lem:profile-strict}
If $x\le y$, then $\prof(x)\preceq\prof(y)$. Under this assumption,
\begin{equation}\label{eq:strict-profile}
 \prof(x)\ne\prof(y)
 \quad\Longleftrightarrow\quad
 \exists j\bigl(\alpha_j<\beta_j\text{ and }\tail(x)\le\beta_j\bigr).
\end{equation}
\end{lemma}
\begin{proof}
The weak comparison follows from the identity matching. If the threshold
increases from $e$ to $f>e$, infinitely many entries of $y$ are at least $e$,
but only finitely many entries of $x$ are; a witness on the right of
\eqref{eq:strict-profile} follows.

If the thresholds agree, any changed or newly created exceptional entry
produces a witness. Conversely, in the absence of such a witness every
coordinate of $y$ at or above the common cut is identical to the corresponding
coordinate of $x$, so the profiles agree. A witness cannot coexist with equal
profiles: the coordinatewise matching of the finite exceptional lists would
then strictly increase at least one summand or introduce a positive new
summand, strictly increasing their finite natural sum.
\end{proof}

\begin{lemma}\label[lemma]{lem:profile-wf}
The profile order is well-founded and set-like: the predecessors of any
profile form a set.
\end{lemma}
\begin{proof}
In an infinite descending chain the cuts form a nonincreasing sequence of
ordinals and hence eventually stabilize. After stabilization, strict descent
of profiles strictly decreases the finite natural sum of their exceptional
entries. This is impossible.

For set-likeness, fix $(e;A)$ and choose an ordinal $M$ at least $e$ and
all elements of $A$. In every predecessor, the cut is at most $e$; heads
below $e$ are bounded by $M$, and heads at or above $e$ must match elements
of $A$ and so are bounded by $M$ as well. All predecessors belong to a set
of finite multisets over $M+1$, together with cuts at most $e$.
\end{proof}

\subsection{Minimality is exactly a rank}
Let $\rho(e;A)$ denote the rank in the profile order, with minimal rank zero.

\begin{theorem}[Rank characterization]\label{thm:rank}
For every sequence $x$,
\begin{equation}\label{eq:rankN}
 \OpN(x)=\rho(\prof(x)).
\end{equation}
Equivalently, $\OpN$ is the rank for the strict relation
\[
 x\eord y\quad\Longleftrightarrow\quad
 x\le y\text{ and the witness in \eqref{eq:strict} exists}.
\]
In particular, $\OpN$ is invariant under permutations and insertion of any
number of zero entries, and is monotone under injective domination.
\end{theorem}
\begin{proof}
By \cref{lem:profile-strict}, the function $x\mapsto\rho(\prof(x))$ satisfies
both admissibility conditions. It remains to show that every admissible
$F$ dominates it.

Proceed by induction on the profile $P=\prof(x)$. For any $Q\prec P$, choose
a sequence $z$ realizing $Q$, using \cref{lem:cutfacts}. By
\cref{lem:matching}, embed its entries injectively into those of $x$. On the
image of the injection put the corresponding entries of $z$; on all unused
coordinates put zero. The resulting sequence $w$ satisfies $w\le x$ and
$\prof(w)=Q$. Insertion of zeros preserves the cut and the exceptional
multiset, directly from the definitions. Since the profiles differ,
\cref{lem:profile-strict} gives $w\eord x$. Therefore
\[
 F(x)\ge F(w)+1\ge\rho(Q)+1.
\]
Taking the supremum over $Q\prec P$ gives $F(x)\ge\rho(P)$.

For the equivalent sequence-rank statement, the relation $\eord$ is
transitive, is well-founded by the same stabilization argument, and has
set-sized predecessor cones. Its rank is weakly monotone: if $x\le y$, every
$\eord$-predecessor of $x$ is one of $y$. Its rank is strictly monotone for
$\eord$, and induction shows that it is dominated by every admissible
operation. It must therefore be the same least operation.

Finally, permutations and zero insertions do not change the profile, and
injective domination is profile comparison by \cref{lem:matching}.
\end{proof}

The quotient of countable ordinal multisets by mutual injective domination
is thus represented exactly by finite profiles. Its rank is the invariant
that \cref{thm:main} evaluates. No comparison of infinite coordinate lists is
needed once the profile is available.

\section{The arithmetic absorption lemma}\label{sec:absorption}

\subsection{Base values and exceptional weights}
Define $\FunB(e)$ for \emph{every} positive ordinal $e$, including cuts that
are not realizable. Define $r_e(a)$ for $a\ge e$ by
\begin{equation}\label{eq:baseweights}
\begin{array}{c|c|c}
\text{type of }e&\FunB(e)&r_e(a)\\ \hline
 e=z+1 &\FunC(z)&\dd{z}{a}\\[2pt]
 e=p+\omega^\xi,\ \xi\text{ successor}
       &\FunA(e)&1+\dd{e}{a}\\[2pt]
 e=p+\omega^\xi,\ \xi\text{ nonzero limit}
       &\FunA(p)\ns\omega^\xi&\dd{p}{a}
\end{array}
\end{equation}
The proposed formula is simply
\begin{equation}\label{eq:U}
 \FunU(e;A)=\FunB(e)\ns\bigns_{a\in A}r_e(a).
\end{equation}
As throughout, repetitions in $A$ are counted.

Each weight is positive and strictly increasing in $a$. Moreover
\begin{equation}\label{eq:weight-cocycle}
 r_e(a)=r_e(t)+\dd{t}{a}\le r_e(t)\ns\dd{t}{a}
 \qquad(e\le t\le a),
\end{equation}
by associativity of ordinary addition and \cref{lem:arith}. The extra leading
one in the middle case of \eqref{eq:baseweights} is compatible with this
identity.

The following convenient estimates follow directly from the definitions:
\begin{equation}\label{eq:basebounds}
 \FunB(e)\le\FunA(e),\qquad
 \FunC(t)\le\FunA(t)\ns\omega.
\end{equation}
For the first bound in the successor case, use
$\FunC(z)\le\FunA(z)\ns\omega=\FunA(z+1)$.
In the final limit case, use
$\FunA(e)=\FunA(p)+\omega^{\xi+1}$.

\subsection{The central estimate}
\begin{lemma}[Absorption]\label[lemma]{lem:absorption}
For $0<e\le t$ and every finite $m$,
\begin{equation}\label{eq:abs1}
 \FunB(e)\ns\bigl(r_e(t)\nprod m\bigr)<\FunC(t).
\end{equation}
If $t$ is a nonzero limit and $e<t$, then the stronger bound holds:
\begin{equation}\label{eq:abs2}
 \FunB(e)\ns\bigl(r_e(t)\nprod m\bigr)<\FunA(t).
\end{equation}
\end{lemma}
\begin{proof}
We prove both assertions simultaneously by transfinite induction on $t$.
The multiplier $m$ is finite, and multiplication by it in these expressions
means a repeated finite natural sum.

\paragraph{Finite $t$.}
Let $t=n\ge1$. Then $e=k+1$ for $k<n$, $\FunB(e)=\omega k$, and
$r_e(n)=n-k$. Hence
\[
 \FunB(e)\ns(r_e(n)\nprod m)
   =\omega k+m(n-k)<\omega(k+1)\le\omega n=\FunC(n).
\]
There is no limit assertion in this case.

\paragraph{A nonzero finite tail.}
Let $t=\lambda+n$, where $\lambda$ is a nonzero limit and $0<n<\omega$.
First suppose $e\le\lambda$. By \eqref{eq:weight-cocycle},
$r_e(t)=r_e(\lambda)+n$, so the left side of \eqref{eq:abs1} is at most
\[
 \bigl(\FunB(e)\ns(r_e(\lambda)\nprod m)\bigr)\ns mn.
\]
The parenthesized ordinal is smaller than $\FunC(\lambda)$ by induction.
Since $\FunC(\lambda)$ is a limit, adding a finite ordinal keeps it smaller.
This proves the required bound, since $\FunC(\lambda)\le\FunC(t)$.

If $\lambda<e\le t$, write $e=\lambda+k+1$ with $k<n$.
Then $\FunB(e)=\FunC(\lambda)+\omega k$ and $r_e(t)=n-k$. The left side is
strictly smaller than
\[
 \FunC(\lambda)+\omega(k+1)
 \le\FunC(\lambda)+\omega n=\FunC(t).
\]

\paragraph{A limit $t$ with $e=t$.}
Write $t=p+\omega^\xi$ as in \eqref{eq:lastmonomial}.
If $\xi$ is a successor, then $\FunB(t)=\FunA(t)$ and $r_t(t)=1$, giving
\[
 \FunB(t)\ns(r_t(t)\nprod m)
  =\FunA(t)+m<\FunA(t)+\omega=\FunC(t).
\]
If $\xi$ is a nonzero limit, then $\FunC(t)=\FunA(t)$ and
\[
 \FunB(t)\ns(r_t(t)\nprod m)
   =\FunA(p)+\omega^\xi(m+1)
   <\FunA(p)+\omega^{\xi+1}=\FunA(t).
\]

\paragraph{A limit $t$ with $e<t$: cuts at or below $p$.}
Assume $e\le p$, so $p>0$, and put $c=r_e(p)$. Then
$r_e(t)=c+\omega^\xi\le c\ns\omega^\xi$. By induction,
\[
 P:=\FunB(e)\ns(c\nprod m)<\FunC(p)\le\FunA(p)\ns\omega.
\]
Thus the expression to be bounded is at most
\[
 P\ns(\omega^\xi\nprod m)
 \le\FunA(p)\ns\omega\ns(\omega^\xi\nprod m)
 <\FunA(p)+\omega^{\xi+1}=\FunA(t).
\]
The last strict inequality follows because every exponent in $\FunA(p)$
is at least $\xi+1$, and $\omega\ns(\omega^\xi\nprod m)$ is smaller than
$\omega^{\xi+1}$ for $\xi>0$.

\paragraph{A limit $t$ with $p<e<t$.}
Write $e=p+\delta$, where $0<\delta<\omega^\xi$.
The pivot used in the remainder defining $r_e$ is at least $p$ and smaller
than $t$. Indeed it is $e-1$, $e$, or the ordinal obtained by deleting one
last monomial from $e$, respectively. Each pivot plus $\omega^\xi$ equals
$t$. Any leading one in the weight is absorbed. Therefore
$r_e(t)=\omega^\xi$.

By \eqref{eq:basebounds} and separation of the CNF exponents,
\[
 \FunB(e)\le\FunA(e)=\FunA(p)\ns\FunA(\delta).
\]
Now $\FunA(\delta)<\omega^{\xi+1}$, so
\[
 \FunA(\delta)\ns(\omega^\xi\nprod m)<\omega^{\xi+1}.
\]
The required expression is again smaller than
$\FunA(p)+\omega^{\xi+1}=\FunA(t)$.
This proves the stronger bound whenever $e<t$ and $t$ is a limit, and
completes the induction.
\end{proof}

The proof did not assume realizability of intermediate cuts. This is useful:
the auxiliary $p$ can have arbitrary cofinality even when $t$ is realizable.

\section{Upper bound: the formula is admissible}\label{sec:upper}

\begin{theorem}\label{thm:upper}
The function $\FunU$ is strictly increasing on strict profile comparisons.
Consequently $x\mapsto\FunU(\prof(x))$ is admissible and
\begin{equation}\label{eq:NleU}
 \OpN(x)\le\FunU(\prof(x)).
\end{equation}
\end{theorem}
\begin{proof}
Suppose $(e;A)\preceq(f;B)$. Match each head $a\in A$ with $a\ge f$ to a
head $b\in B$ with $a\le b$. All heads of $A$ below $f$ are unmatched.

\paragraph{Equal cuts.}
If $e=f$, every head in $A$ is matched. Positivity and strict monotonicity of
$r_e$ show $\FunU(e;A)\le\FunU(e;B)$. If the profiles differ, at least one
matched value increases or an additional positive weight occurs, so the
inequality is strict.

It remains to assume $e<f$. Put $m=|A|$.

\paragraph{The upper cut is a successor.}
Let $f=t+1$. Then $e\le t$. Every unmatched $a$ satisfies $a\le t$, so
$r_e(a)\le r_e(t)$. For a matched $a$, \eqref{eq:weight-cocycle} gives
\[
 r_e(a)\le r_e(t)\ns\dd{t}{a}
          =r_e(t)\ns r_f(a).
\]
Consequently, writing $A_{\ge f}$ for the matched part of $A$,
\begin{align*}
 \FunU(e;A)
 &\le\bigl(\FunB(e)\ns(r_e(t)\nprod m)\bigr)
                       \ns\bigns_{a\in A_{\ge f}}r_f(a)\\
 &<\FunC(t)\ns\bigns_{a\in A_{\ge f}}r_f(a)
 \le\FunB(f)\ns\bigns_{b\in B}r_f(b)=\FunU(f;B).
\end{align*}
The strict step is \cref{lem:absorption}.

\paragraph{The upper cut is a limit with successor last exponent.}
Write $f=p+\omega^\xi$ with $\xi$ a successor.
Unmatched weights are at most $r_e(f)$. For matched heads,
\[
 r_e(a)\le r_e(f)\ns\dd{f}{a}
          \le r_e(f)\ns r_f(a).
\]
Use \eqref{eq:abs2} with $t=f$ to replace the parenthesized bulk term by an
ordinal strictly below $\FunA(f)=\FunB(f)$. Matching the remaining weights
into those of $B$ then proves strictness just as above.

\paragraph{The upper cut is a limit with limit last exponent: $e\le p$.}
Now $f=p+\omega^\xi$, with $\xi$ a nonzero limit, and
$\FunB(f)=\FunA(p)+\omega^\xi$. Suppose first that $e\le p$.
Set $c=r_e(p)$. Heads below $p$ have weight at most $c$.
For $p\le a<f$, the weight is at most $c\ns\dd{p}{a}$, with
$\dd{p}{a}<\omega^\xi$. For a matched $a\ge f$, the weight is at most
$c\ns r_f(a)$. A finite natural sum of the extra remainders below $f$ is
some $\eta<\omega^\xi$. Therefore
\[
 \FunU(e;A)\le
 \bigl(\FunB(e)\ns(c\nprod m)\ns\eta\bigr)
      \ns\bigns_{a\in A_{\ge f}}r_f(a).
\]
By \eqref{eq:abs1}, the bulk before adding $\eta$ is smaller than
$\FunC(p)\le\FunA(p)\ns\omega$. Since $\xi$ is a nonzero limit,
$\omega\ns\eta<\omega^\xi$. The full bulk is thus strictly smaller than
$\FunA(p)+\omega^\xi=\FunB(f)$. Matching the remaining weights proves the
claim. When $p=0$, this subcase is empty.

\paragraph{The upper cut is a limit with limit last exponent: $p<e<f$.}
Write $e=p+\delta$, with $0<\delta<\omega^\xi$. The limit-exponent clause
of \cref{lem:arith} gives
\[
 \FunB(e)\le\FunA(e)=\FunA(p)\ns\FunA(\delta),
 \qquad\FunA(\delta)<\omega^\xi.
\]
The pivot for every remainder defining $r_e$ lies in $[p,f)$.
It follows that $r_e(a)<\omega^\xi$ for unmatched heads $a<f$.
For a matched head $a\ge f$, that pivot plus $\omega^\xi$ is $f$, so
\[
 r_e(a)=\omega^\xi+\dd{f}{a}=\dd{p}{a}=r_f(a).
\]
In the case with a leading one, that one is absorbed by $\omega^\xi$.
The base and all unmatched weights together are therefore at most
$\FunA(p)\ns\eta$ for some $\eta<\omega^\xi$, strictly below $\FunB(f)$.
Matching the remaining weights finishes the proof.

We have proved strictness whenever the profiles differ and weak monotonicity
otherwise. Apply \cref{lem:profile-strict} and the defining minimality of
$\OpN$, or equivalently \cref{thm:rank}, to obtain \eqref{eq:NleU}.
\end{proof}

\section{Lower bound and completion of the proof}\label{sec:lower}

\subsection{The published input, stated precisely}
The weaker operation in Lipparini's article is denoted there by $\#^S$;
we call it $\OpS$. Define $\zeta(x)=z$ when $\tail(x)=z+1$, and
$\zeta(x)=\tail(x)$ when the cut is a limit. Its strictness requirement is
\eqref{eq:strict} with $\tail(x)\le\beta_j$ replaced by
$\zeta(x)<\beta_j$.

\begin{proposition}[Input from Lipparini]\label[proposition]{prop:importedS}
The least weakly monotone operation with that weaker strictness requirement
is
\begin{equation}\label{eq:Sformula}
 \OpS(x)=
 \begin{cases}
 \FunA(z)\ns\displaystyle\bigns_{i\in E}\dd{z}{\alpha_i},
      &e=z+1,\\[5pt]
 \FunA(e)\ns\displaystyle\bigns_{i\in E}\dd{e}{\alpha_i},
      &e=p+\omega^\xi,\ \xi\text{ successor},\\[5pt]
 \FunA(p)\ns\omega^\xi\ns\displaystyle\bigns_{i\in E}\dd{p}{\alpha_i},
      &e=p+\omega^\xi,\ \xi\text{ nonzero limit}.
 \end{cases}
\end{equation}
In particular $\OpN(x)\ge\OpS(x)$.
\end{proposition}
\begin{proof}[Source and deduction]
The formula and minimality are Definition 3.1 and Theorem 3.4 of
\cite{Lipparini2026extended}. The last inequality follows because every
admissible operation here satisfies the weaker strictness condition:
if $\zeta(x)<\beta_j$, then $\tail(x)\le\beta_j$, in both the successor and
limit cases. We use the cited minimality theorem, not merely the displayed
formula as an experimentally plausible lower bound.
\end{proof}

This is the only external theorem about infinitary ordinal operations used
in the lower-bound argument. The product-rank lemma and all new threshold
estimates have been proved above.

\subsection{A finite-head lower bound}
Write $\cseq{z}=(z,z,z,\ldots)$.

\begin{lemma}\label[lemma]{lem:finite-head}
For $a_1,\ldots,a_h\ge z$, put $\delta_j=\dd{z}{a_j}$. Then
\begin{equation}\label{eq:finite-head}
 \OpN(a_1,\ldots,a_h,z,z,\ldots)
 \ge\OpN(\cseq{z})\ns\delta_1\ns\cdots\ns\delta_h.
\end{equation}
If $\tail(x)=z+1$, the same bound holds using the exceptional entries of $x$.
\end{lemma}
\begin{proof}
Let $Q$ be the set of sequences coordinatewise below $\cseq{z}$, with the
strict relation $\eord$. Its rank at $\cseq{z}$ is $\OpN(\cseq{z})$ by
\cref{thm:rank}: this set contains the entire predecessor cone.
For $b\in Q$ and $0\le\rho_j\le\delta_j$, map
\[
 (b,\rho_1,\ldots,\rho_h)
 \longmapsto (z+\rho_1,\ldots,z+\rho_h,b_0,b_1,\ldots).
\]
Adding finitely many heads does not change the tail cut. A strict tail move
remains an $e$-strict move. A strict increase of any $\rho_j$ makes the target
head larger than $z$, hence at least $z+1$, while the lower tail cut is at
most $z+1$. It too is $e$-strict. The map therefore preserves strict order
from the product of $Q$ and the ordinal chains $\delta_j+1$.
Apply \cref{lem:product-rank} at
$(\cseq{z},\delta_1,\ldots,\delta_h)$.

For a general $x$ with cut $z+1$, keep its exceptional entries and an
infinite subsequence of entries equal to $z$, and replace all other entries
by zero. This gives a sequence below $x$ whose profile is the same as that
of the displayed finite-head sequence. Apply \cref{thm:rank} and weak
monotonicity.
\end{proof}

\subsection{The finite correction at a limit cut}
\begin{lemma}\label[lemma]{lem:limit-lower}
Suppose $e=\tail(x)=p+\omega^\xi$ with $\xi$ a successor. Set
\begin{equation}\label{eq:k}
 k(x)=|\{i\in E(x):\alpha_i<e+\omega\}|.
\end{equation}
Then
\begin{equation}\label{eq:limitlower}
 \OpN(x)\ge\OpS(x)+k(x)=\FunU(\prof(x)).
\end{equation}
\end{lemma}
\begin{proof}
Let $J$ be the set counted by $k(x)$. For $i\in J$, write
$\alpha_i=e+n_i$ with $n_i<\omega$, and put $m=\sum_{i\in J}n_i$.
Replace all these entries by zero, obtaining $x_0$. Its cut remains $e$.
The middle case of \eqref{eq:Sformula} gives
\[
 \OpS(x)=\OpS(x_0)+m.
\]
Here a finite natural summand is the same as ordinary addition on the right,
so the equality remains valid even when $\OpS(x_0)$ has a finite part.

Restore the $i$th modified coordinate through
\[
 0\longrightarrow e\longrightarrow e+1\longrightarrow\cdots
                         \longrightarrow e+n_i.
\]
There are $n_i+1$ strict steps, each satisfying \eqref{eq:strict}. The cut
stays $e$ throughout. Restoring all modified coordinates therefore gives
\[
 \OpN(x)\ge\OpN(x_0)+(m+k(x))
          \ge\OpS(x_0)+(m+k(x))=\OpS(x)+k(x).
\]
Finally, $1+d$ is $d+1$ for finite $d$, and equals $d$ for every infinite
ordinal $d$. Applying this fact to the weights in \eqref{eq:baseweights}
shows that $\FunU=\OpS+k(x)$ in precisely this case.
\end{proof}

\subsection{Constant sequences and the extra \texorpdfstring{$\omega$}{omega}}
\begin{lemma}\label[lemma]{lem:constant-lower}
For every ordinal $t$,
\begin{equation}\label{eq:constantlower}
 \OpN(\cseq{t})\ge\FunC(t).
\end{equation}
\end{lemma}
\begin{proof}
If $\chi(t)=0$, the first case of \eqref{eq:Sformula} and
\cref{prop:importedS} give
$\OpN(\cseq{t})\ge\FunA(t)=\FunC(t)$.

Suppose $\chi(t)=1$. Write $t=\lambda+n$ with $n<\omega$, where the last
exponent of the nonzero limit $\lambda$ is a successor. Such a $\lambda$
has cofinality $\omega$: writing its final monomial as
$\omega^{\eta+1}=\omega^\eta\cdot\omega$ supplies a countable cofinal tail,
also when $\eta$ is uncountable.

Choose a sequence cofinal below $\lambda$. For each finite $m$, append
$m$ copies of $\lambda$. The cut is $\lambda$; its weaker value is
$\FunA(\lambda)$, and \cref{lem:limit-lower} gives lower bound
$\FunA(\lambda)+m$. All entries are at most $\lambda$, so monotonicity yields
\[
 \OpN(\cseq{\lambda})\ge
 \sup_{m<\omega}\bigl(\FunA(\lambda)+m\bigr)
 =\FunA(\lambda)+\omega=\FunC(\lambda).
\]

Proceed by induction on the finite $n>0$. Append $m$ copies of
$\lambda+n$ to the constant $\lambda+(n-1)$ tail. The remainders of these
heads over the tail value are all one. By \cref{lem:finite-head}, its value
is at least
\[
 \OpN(\cseq{\lambda+(n-1)})\ns m
 \ge\FunC(\lambda+(n-1))\ns m.
\]
It is dominated by $\cseq{\lambda+n}$. Taking the supremum over $m$ and
using \eqref{eq:Cfinite} gives
\[
 \OpN(\cseq{\lambda+n})\ge
 \FunC(\lambda+(n-1))+\omega=\FunC(\lambda+n).
\]
\end{proof}

\subsection{Completion and a dependency check}
\begin{proof}[Proof of \cref{thm:main}]
The upper bound is \cref{thm:upper}. For the reverse inequality, if the cut
is a limit with successor last exponent, apply \cref{lem:limit-lower}.
If it is a limit with nonzero limit last exponent, $\FunU=\OpS$, so
\cref{prop:importedS} suffices.

If the cut is $z+1$, combine \cref{lem:finite-head,lem:constant-lower}:
\[
 \OpN(x)\ge\OpN(\cseq{z})\ns\bigns_{i\in E(x)}\dd{z}{\alpha_i}
 \ge\FunC(z)\ns\bigns_{i\in E(x)}\dd{z}{\alpha_i}
 =\FunU(\prof(x)).
\]
Thus both bounds agree in every case.
\end{proof}

There is no circular use of the constant-sequence formula. The order of
proof is: the independent lower bound $\OpN\ge\OpS$; finite-product ranks;
the finite correction at a limit cut; the constant lower bound; and finally
the general successor-cut lower bound. The upper-bound argument was proved
without assuming any of these lower-bound conclusions.

\section{Consequences and worked examples}\label{sec:consequences}

\subsection{Exactly how the stronger sum differs}
\begin{corollary}[Complete correction rule]\label[corollary]{cor:correction}
With the previous notation,
\begin{equation}\label{eq:correction}
 \OpN(x)=
 \begin{cases}
 \OpS(x)\ns\bigl(\omega\cdot\chi(z)\bigr),&e=z+1,\\
 \OpS(x)+k(x),&e\text{ limit with successor last exponent},\\
 \OpS(x),&e\text{ limit with nonzero limit last exponent}.
 \end{cases}
\end{equation}
In particular,
\begin{equation}\label{eq:smallgap}
 \OpS(x)\le\OpN(x)\le\OpS(x)\ns\omega.
\end{equation}
The upper bound is attained exactly in the successor-cut case with
$\chi(z)=1$. The two operations have identical CNF terms of exponent at least
$2$.
\end{corollary}
\begin{proof}
Compare \cref{thm:main} with \eqref{eq:Sformula}. The middle case was computed
in \cref{lem:limit-lower}. A finite increment is strictly smaller than adding
one natural summand $\omega$, and natural addition by $\omega$ is strictly
increasing. Only the coefficients of $1$ or $\omega$ can change.
\end{proof}

Thus equality $\OpN(x)=\OpS(x)$ holds exactly when $\chi(z)=0$ in the
successor case, when $k(x)=0$ in the middle case, or always in the final case.
Both the arbitrary finite corrections and the full natural $\omega$
correction occur. In particular, the stronger operation is not obtained by
blindly adding one to every exceptional remainder.

\subsection{Constant values and bounded-multiset height}
\begin{corollary}\label[corollary]{cor:constant}
For every ordinal $t$,
\begin{equation}\label{eq:constant}
 \OpN(t,t,t,\ldots)=\FunC(t)
 = (t\nprod\omega)\ns\bigl(\omega\cdot\chi(t)\bigr).
\end{equation}
Let $\mathscr M_t$ be the poset of countably infinite multisets with entries
in $t+1$, quotiented by mutual injective domination. With height defined as
$\sup_p(\rk(p)+1)$,
\begin{equation}\label{eq:height}
 \operatorname{ht}(\mathscr M_t)=\FunC(t)+1.
\end{equation}
\end{corollary}
\begin{proof}
The constant sequence has cut $t+1$ and no exceptional entries, proving
\eqref{eq:constant}. It is the greatest element of the bounded domination
poset: every sequence with entries at most $t$ injects into it. Its
predecessor cone is the whole poset, so its rank is the unrestricted rank
$\FunC(t)$. The stated height convention adds one.
\end{proof}

For instance, these heights are $\omega m+1$ for positive finite $t=m$,
$\omega^2+\omega+1$ for $t=\omega$, and $\omega^{\omega+1}+1$ for
$t=\omega^\omega$. The additive distinction at the least limit exponent
therefore has a direct order-theoretic interpretation.

\subsection{A table of profiles}
A row described as a cofinal tail contains only entries below the named cut,
apart from the explicitly listed exceptional heads. Finite changes below
the cut do not affect the answer.

\begin{center}
\small
\begin{tabular}{@{}p{.45\textwidth}p{.23\textwidth}p{.23\textwidth}@{}}
\toprule
Sequence or profile & $\OpS$ & $\OpN$\\
\midrule
Constant $0$ & $0$ & $0$\\[3pt]
Head $5$, otherwise $0$ & $5$ & $5$\\[3pt]
Cofinal natural-number tail & $\omega^2$ & $\omega^2$\\[3pt]
One head $\omega$ over that tail & $\omega^2$ & $\omega^2+1$\\[3pt]
Two heads $\omega$ over that tail & $\omega^2$ & $\omega^2+2$\\[3pt]
Head $\omega+5$ over that tail & $\omega^2+5$ & $\omega^2+6$\\[3pt]
Head $\omega\cdot2+3$ over that tail
 & $\omega^2+\omega+3$ & $\omega^2+\omega+3$\\[3pt]
Constant $\omega$ & $\omega^2$ & $\omega^2+\omega$\\[3pt]
Constant $\omega+1$ & $\omega^2+\omega$ & $\omega^2+\omega\cdot2$\\[3pt]
Cofinal tail below $\omega^\omega$ & $\omega^\omega$ & $\omega^\omega$\\[3pt]
One head $\omega^\omega$ over that tail
 & $\omega^\omega\cdot2$ & $\omega^\omega\cdot2$\\[3pt]
Constant $\omega^\omega$ & $\omega^{\omega+1}$ & $\omega^{\omega+1}$\\[3pt]
Constant $\omega^{\omega+1}$ & $\omega^{\omega+2}$ & $\omega^{\omega+2}+\omega$\\
\bottomrule
\end{tabular}
\end{center}

The basic $\omega$ examples agree with those in
\cite[Remark 6.1(c)]{Lipparini2026extended}; the general formula explains the
same phenomenon at arbitrary CNF exponents.

\begin{example}[Why the direction of $1+d$ matters]\label[example]{ex:leftone}
Let the cut be $\omega$, with the single exceptional head
$a=\omega\cdot2+3$. Its remainder is $\dd{\omega}{a}=\omega+3$.
The correct weight is
\[
 1+(\omega+3)=\omega+3,
\]
not $(\omega+3)+1=\omega+4$. Thus the stronger and weaker operations agree
on this input. A rule using $d+1$ would falsely predict an additional one.
\end{example}

\begin{example}[A nontrivial prefix in the limit-exponent case]
Let $e=\omega^{\omega+1}+\omega^\omega\cdot2$, and suppose the exceptional
heads are $a_1=e$ and $a_2=e+7$. Deleting one final monomial gives
$p=\omega^{\omega+1}+\omega^\omega$, not $\omega^{\omega+1}$. Then
\[
 \dd{p}{a_1}=\omega^\omega,\qquad
 \dd{p}{a_2}=\omega^\omega+7,
\]
and
\[
 \OpN(x)=\OpS(x)=
 \omega^{\omega+2}+\omega^{\omega+1}+\omega^\omega\cdot3+7.
\]
This illustrates why multiplicity in the final CNF term must be handled
before evaluating the formula.
\end{example}

\subsection{A structural explanation of the correction}
At a limit cut whose last exponent is a successor, an exceptional entry
exactly equal to the cut has zero weaker remainder. But moving a coordinate
from below the cut to the cut is a strict move for $\OpN$. This accounts for
the leading one in $1+d$; only finite remainders retain that one.

A constant limit tail can dominate sequences with arbitrarily many such
finite correction units. Their supremum forces one extra $\omega$. That
extra term persists along the subsequent finite successor tail of the
constant value. When the least exponent is itself a limit, the lower-cut
base ranks already grow cofinally through the relevant $\omega^\xi$ block;
the absorption estimates show that no extra low-order term is required.
This explanation is reflected quantitatively in \cref{lem:absorption} and
\cref{lem:constant-lower}, rather than being used as a substitute for them.

\section{Finite caps: an independent rank calculation}\label{sec:finitecaps}

A finite test can check the order-theoretic modeling independently of the
transfinite arithmetic. It also demonstrates why finite testing cannot prove
an ordinal-rank formula by taking a naive supremum.

Fix integers $M,H\ge0$. Consider profiles with entries in $\{0,\ldots,M\}$
and at most $H$ exceptional heads. Every such profile has the form
$(z+1;A)$ with $0\le z\le M$ and $A$ a multiset of integers strictly larger
than $z$. This is a finite induced subposet of the full profile order.

\begin{theorem}[Finite-cap rank]\label{thm:finitecap}
The rank in this finite induced subposet is
\begin{equation}\label{eq:finitecap}
 \rho_H(z+1;A)=z(H+1)+\sum_{a\in A}(a-z).
\end{equation}
In particular, the constant-$M$ profile has finite-cap rank $M(H+1)$.
\end{theorem}
\begin{proof}
Let the displayed expression be $R$. For equal cuts, positivity and
monotonicity of the head weights show that $R$ strictly increases on strict
comparison. Suppose the lower and upper tail values are $z<t$.
Each lower head contributes at most $t-z$ plus its upper-cut remainder
when it remains exceptional; a swallowed head contributes at most $t-z$.
There are at most $H$ heads. Thus for $(z+1;A)\preceq(t+1;B)$,
\begin{align*}
 R(z+1;A)
 &\le z(H+1)+H(t-z)+\sum_{b\in B}(b-t)\\
 &=R(t+1;B)-(t-z)<R(t+1;B).
\end{align*}
Hence the actual finite rank is at most $R$.

For the reverse bound, induct on $z$. The constant-zero profile has rank
zero. The profile with tail value $z-1$ and $H$ heads equal to $z$ lies below
the constant-$z$ profile. By the inductive lower bound, its rank is at least
\[
 (z-1)(H+1)+H=z(H+1)-1.
\]
So the constant-$z$ profile has rank at least $z(H+1)$.
From it, create each desired head at $z+1$ and increase it one unit at a time
to its final value. This is a chain of length $\sum_{a\in A}(a-z)$, using at
most $H$ heads. Therefore the rank of $(z+1;A)$ is at least $R$.
\end{proof}

\begin{remark}[Ranks do not commute with these increasing unions]
For fixed finite $M\ge2$, the finite-cap subposets increase with $H$ and their
union is the entire bounded profile poset. Nevertheless
\[
 \sup_{H<\omega}M(H+1)=\omega
 \quad\text{whereas}\quad
 \rho(M+1;\emptyheads)=\omega M>\omega.
\]
There is no contradiction. The finite-cap subposets are not closed under
all predecessors in the full poset, and rank is not continuous under these
arbitrary increasing unions. This is a concrete reason not to mistake
successful finite-rank experiments for a transfinite minimality proof.
\end{remark}

\section{Implementation and verification}\label{sec:implementation}

\subsection{An exact, finite-notation implementation}
The accompanying \texttt{code/ordinals.py} represents ordinals below
$\epsilon_0$ by finite hereditary Cantor normal forms. A normal form is a
finite decreasing tuple of pairs consisting of a recursively represented
exponent and a positive integer coefficient. Zero is the empty tuple.
There is no floating-point arithmetic and no truncation of infinite series.

The implemented operations are comparison, ordinary addition, natural
addition, ordinary-addition remainder, finite natural multiplication,
natural multiplication by $\omega$, extraction of the last monomial,
$\chi$, profile comparison, $\FunU$, and the weaker $\OpS$ formula.
The code distinguishes the two additions in the middle branch explicitly:
\begin{lstlisting}[language=Python]
return ONE + e.right_remainder(a)  # ordinary 1+d, not d+1
\end{lstlisting}
The input to evaluation is a validated profile, not an unevaluated infinite
stream. Normal-form size and recursive exponent comparison determine the
cost of the basic arithmetic. No complexity claim is made for arbitrary
ordinal notation systems, and the implementation does not represent
$\epsilon_0$ itself or uncountable ordinals.

The mathematics of \cref{thm:main} is not restricted to this implementation's
domain. In particular, the proof of the absorption lemma allowed arbitrary
ordinal exponents and arbitrary intermediate cofinalities.

\subsection{What was actually checked}
The deterministic test run used seed \texttt{20260919}, a basis of 223 exact
ordinals, 10,000 randomized cases, and an exhaustively compared family of
1,000 small profiles. The basis includes limit and successor exponents,
multiple CNF terms, nonunit coefficients, and finite tails. The results are
stored in \texttt{checks/results.json} and \texttt{checks/run.txt}.

\begin{center}
\begin{tabular}{@{}lr@{}}
\toprule
Category & Counted checks\\
\midrule
Arithmetic boundary, pair, and triple checks & 120,127\\
Remainder and remainder-cocycle checks & 34,976\\
General absorption inequality & 222,777\\
Stronger limit absorption inequality & 63,792\\
Correction formula on generated profiles & 10,000\\
Generated strict or equal comparable profiles & 10,000\\
Exhaustive comparable pairs of small profiles & 272,272\\
Independent finite-poset rank computations & 912\\
Explicit regression examples & 13\\
\midrule
Total & 734,869\\
\bottomrule
\end{tabular}
\end{center}

Every counted test passed. Some counted cases contain several assertions;
these figures follow the categories reported by the test script, not a claim
that every Python assertion is counted separately. The finite-poset ranks
were computed from predecessor sets without using \eqref{eq:finitecap} to
compute the ranks; the formula was checked afterward. The test domains were
$1\le M\le5$ and $0\le H\le4$.

These checks are designed to expose formula errors, wrong remainders,
incorrect additions, and broken profile comparisons. They do not prove
well-foundedness for arbitrary ordinals, justify a proper-class operation,
or establish minimality. Those claims rest on the preceding proofs and
the identified published input. No formal verification in Lean, Isabelle,
or another proof assistant was performed.

\subsection{Why arbitrary cut extraction is not computable}\label{subsec:computability}
There is an elementary obstruction even for computable sequences with
entries only $0$ and $1$. Given a machine $T$, define $a_n=1$ if $T$ has
halted within $n$ steps, and $a_n=0$ otherwise. The sequence is computable
uniformly from $T$. If $T$ never halts, the cut is $1$ and the value is zero.
If it halts, the cut is $2$ and the value is $\omega$. An algorithm that
always obtained either the exact cut or the exact value from such a program
would decide whether $T$ halts.

Thus an \emph{explicit mathematical formula} must not be confused with a
uniform decision procedure for arbitrary program-presented infinite inputs.
The supplied evaluator correctly accepts finite profile data; it makes no
promise to discover such data from an arbitrary stream.

\section{Scope, audit points, and further questions}\label{sec:scope}

\subsection{What is established by the argument}
Relative to Lipparini's proved minimality theorem for $\OpS$, the article
proves a formula for $\OpN$ on \emph{all} countable sequences of arbitrary
ordinals, not merely on eventually constant sequences or ordinals below
$\epsilon_0$. The formula is pointwise least among the operations satisfying
the exact two admissibility axioms. The correction rule, constant values,
profile-rank interpretation, bounded-multiset heights, and finite-cap ranks
are deductions, not additional conjectures.

The countable index set is essential to the stated setting. At a limit cut
it forces cofinality $\omega$; changing the index cardinal changes both the
possible profiles and the finite-exception criterion. No uncountable-arity
generalization, cardinal-consistency result, or claim about large cardinals
is asserted here.

\subsection{Foundations and the location of the proof burden}
The profile rank is defined on set-sized predecessor cones. Restricting to
any sufficiently large ordinal bound gives an ordinary set-sized
well-founded relation; these ranks agree when the bound is enlarged because
predecessor cones are unchanged. Thus the notation for a class function
causes no need to minimize over a proper class of arbitrary functions.
The same interpretation applies to the sequence-rank proof and to the
finite-product argument on the bounded set $Q$.

A useful audit focuses on four points: the strict relation uses the lower
sequence's cut; the remainder must solve ordinary $a+d=b$; the middle weight
must be $1+d$ rather than $d+1$; and the final-monomial deletion must remove
one copy, not an entire CNF term. The computational tests explicitly exercise
these distinctions. The most substantial proof obligations are the
cross-cut comparisons in \cref{thm:upper} and the finite-product lower bound
in \cref{lem:finite-head}.

\subsection{Natural next questions, not claimed results}
The rank formula suggests looking for a direct shuffled-order realization of
$\OpN$, analogous in spirit to characterizations of other natural ordinal
operations. A separate question is which restricted forms of associativity
survive after the correction in \eqref{eq:correction}. Neither question is
settled by simply knowing the CNF formula. The present note does not assume
answers to them, and it does not claim that they are new open problems.

The immediate research task is independent verification of the proposed
solution and a fuller priority check. A proof-assistant development could
begin with the set-sized profile order, ordinary ordinal remainders, and the
absorption lemma, leaving general stream computability outside its scope.

\appendix
\section{A compact evaluation recipe}\label{app:recipe}
Given a valid cut $e$ and a finite exceptional multiset $A$:
\begin{enumerate}
\item If $e=z+1$, remove the finite tail of $z$ to compute $\chi(z)$.
      Compute $z\nprod\omega$, add a natural $\omega$ when $\chi(z)=1$,
      and add the finite natural sum of the remainders $\dd{z}{a}$.
\item Otherwise delete one last monomial of $e$ to obtain $e=p+\omega^\xi$.
      If $\xi$ is a successor, compute $e\nprod\omega$ and add the finite
      natural sum of $1+\dd{e}{a}$.
\item If $\xi$ is a nonzero limit, compute
      $(p\nprod\omega)\ns\omega^\xi$ and add the finite natural sum of
      $\dd{p}{a}$.
\end{enumerate}
Only finite CNF operations occur after the profile is known. Zero exceptional
heads are impossible because $e>0$; zero entries in the original sequence
remain harmless.

\section{Reproducing the files and tests}
From the archive root, using Python 3.10 or later:
\begin{lstlisting}
python3 code/demo.py
python3 code/verify.py --output checks/results.json
\end{lstlisting}
No third-party Python packages are required. The default seed and sample
count reproduce all mathematical check counts; elapsed time depends on the
machine. To rebuild the article with a TeX installation containing
\texttt{newpx}, \texttt{tcolorbox}, and the other standard packages used in
the preamble, run:
\begin{lstlisting}
sh build.sh
\end{lstlisting}
The archive includes the editable LaTeX source, the compiled PDF, the exact
ordinal implementation, the deterministic tests, a small demonstration,
the recorded results, and a README explaining scope and research status.
It deliberately contains no downloaded font files or copies of the source
paper.

\begin{thebibliography}{9}
\bibitem{Lipparini2026extended}
Paolo Lipparini.
\newblock \emph{Monotone infinitary operations on ordinals (extended version)}.
\newblock arXiv:2505.00424v2, 30 April 2026, 42 pp.
\newblock \url{https://arxiv.org/abs/2505.00424v2}.
\newblock The numbering used here is that of version 2: Definition 2.4,
Theorem 2.6.a, Definition 3.1, Theorem 3.4, Remark 6.1, and Problem 6.2.

\bibitem{Lipparini2026published}
Paolo Lipparini.
\newblock A monotone infinitary operation on ordinals.
\newblock \emph{Mathematical Logic Quarterly} \textbf{72}, no.~2 (2026),
e70019. First published 24 March 2026.
\newblock \url{https://doi.org/10.1002/malq.70019}.
\newblock The extended arXiv version supplies the proofs and numbering
consulted for this note.
\end{thebibliography}
\end{document}
